\section{Introduction}

Interventions---removing a layer, patching an activation, or injecting
noise---are widely used to infer redundancy, necessity, robustness, or causal
roles in neural networks. Such interpretations are strongest when reasonable
intervention procedures agree and weaker when the response depends on the
measurement protocol.

This concern is established empirically. Language models can preserve many
next-token predictions after layer deletion or swapping
\citep{lad2025remarkable}, while \citet{garcia2026nofreeswap} shows that
output-grounded layer equivalence changes across training and differs between
replacement and interchange. Intervention audits also show dependence on the
comparator and matching procedure \citep{luo2026steercheck}. We therefore ask
a narrower question: \emph{can protocol-dependent responses be explained by
differences in functional predictive damage?}

The explanation is plausible. Nominally equal interventions can change output
distributions differently, while interventions far apart in activation space
can be close in function space. If matching forward KL and target-token NLL
removes protocol dependence, intervention identity adds little beyond a
severity budget. If not, raw size and those output-level measures are
insufficient summaries.

We study independently pretrained 70M- and 160M-parameter \model{} models
\citep{biderman2023pythia}. For a fixed corpus and checkpoint, we intervene on
an interior block and measure the retained fraction $S$ of intact next-token
top-1 predictions; $D_S=1-S$ is top-1 damage and
$\deltas=S_{\mathrm{late}}-S_{\mathrm{early}}$ is the training endpoint change.
This operational metric is not new: it belongs to the relative-accuracy family
\citep{lad2025remarkable} and complements the divergent-token rate on identical
positions \citep{deiseroth2024dtm}.

The block-bypass endpoint is negative in all nine 160M runs, including a
prospectively held-out run, and repeats at 70M and on a HellaSwag-derived
stream. This establishes a stable outcome, not novelty for the broad training
phenomenon. Within residual attenuation, complementary matching shows that
KL/NLL damage is more informative than raw activation displacement, without
establishing KL/NLL causality.

The central test prospectively targets activation-noise strengths, using only
KL and NLL, to match block interventions by run, checkpoint, and layer. After a
tie-handling inconsistency was withdrawn and repaired, the canonical analysis
retains a family residual of $-0.0168$ (run-bootstrap 95\% interval
$[-0.0200,-0.0126]$): matched activation noise produces less top-1 damage than
block attenuation on average. A frozen output-logit and decision-boundary
feature block reduces this residual by 20.2\% but does not eliminate it.

Our contributions are exactly three:
\begin{enumerate}
    \item We establish an independently replicated empirical base for a fixed
    block-bypass response across two small \model{} scales and two English
    evaluation streams, including held-out, confidence, and layer controls.
    \item Complementary magnitude- and damage-matched contrasts show that
    output-level predictive damage is more informative than raw activation
    displacement within the frozen residual-attenuation assay.
    \item Prospective joint KL/NLL targeting and a corrected deterministic
    top-1 analysis show that predictive-damage matching does not eliminate the
    block/noise difference; simple output geometry explains only part of it.
\end{enumerate}

We do not infer universal specialization, large-model scaling, downstream
capability loss, KL/NLL causality, or an identified internal mechanism.
